% =============================================================================
%  Thesis Section — 3.4  Conception des données
%  Design-level description: MCD, MLD, Feature Store, Temporal management
% =============================================================================
%  Prerequisites in preamble:
%    \usepackage{graphicx, booktabs, multirow, amsmath, hyperref}
%    \graphicspath{{ml-backend/simulator/}}
%  Figures referenced:
%    figures/mcd_diagram.png
%    figures/mld_tiers_diagram.png
% =============================================================================


% ─────────────────────────────────────────────────────────────────────────────
% 3.4  Conception des données
% ─────────────────────────────────────────────────────────────────────────────

\section{Conception des données}
\label{sec:data-conception}

The PIOS platform handles four families of data: (i)~operational data
(users, alerts, notifications, ML model registry); (ii)~synthetic datasets
that reproduce the statistical properties of a real blood-supply network;
(iii)~feature views consumed by training and inference pipelines; and
(iv)~temporal series that feed the simulator and the reinforcement-learning
agents.  The subsections below present, in turn, the conceptual model, its
logical translation, the feature-store architecture and the strategy for
managing time-dependent data.


% ─────────────────────────────────────────────────────────────────────────────
% 3.4.1  Modèle Conceptuel (MCD)
% ─────────────────────────────────────────────────────────────────────────────

\subsection{Modèle Conceptuel de Données (MCD)}
\label{sec:mcd}

The conceptual model identifies the domain entities, their principal
attributes and the associations between them, independently of any storage
technology.  Figure~\ref{fig:mcd} shows the entity-relationship diagram;
the paragraphs below motivate the design of each domain group.

\begin{figure}[htbp]
\centering
\includegraphics[width=\linewidth]{figures/mcd_diagram.png}
\caption{Modèle Conceptuel de Données.  Solid arrows: foreign-key
         associations; dashed arrows: soft (UUID-based) references.
         Cardinalities use the $1{:}N$ / $M{:}N$ notation.}
\label{fig:mcd}
\end{figure}

\paragraph{Authentication \& collaboration.}

A central \textbf{User} entity supports email-based authentication and
role management ($M{:}N$ with \textbf{Group} and \textbf{Permission}).
Users interact through \textbf{Conversations} (team, mission or direct
threads), each containing ordered \textbf{Messages} with a priority tag.
Separating conversations by type allows the UI to filter views without
adding structural complexity.

\paragraph{Alerting \& notifications.}

The alerting subsystem follows a \emph{rule $\to$ event $\to$ history}
pattern.  An \textbf{AlertRule} encodes reusable trigger logic—scope
(global, centre, hospital, blood type), trigger type (threshold gap,
stock-and-forecast, stale alert, manual), severity, de-duplication window
and escalation delay.  When a rule fires it creates an
\textbf{AlertEvent} whose status progresses through \emph{open},
\emph{acknowledged}, \emph{escalated} and \emph{resolved}.  Every
transition is appended to an \textbf{AlertEventHistory} record, providing
a full audit trail.  A \textbf{Notification} entity handles delivery
(push / in-app) and is linked from AlertEvent by a soft UUID reference
rather than a foreign key, which decouples alert lifecycle management from
the notification channel.

\paragraph{ML model registry \& observability.}

Every deployed model is represented by an \textbf{MLModelConfig} entry
that records its type, feature schema, example payloads and default
inputs.  Inference calls produce \textbf{PredictionLog} records tied to
the calling User and the config used, enabling usage analytics and latency
monitoring.  An \textbf{OrchestratorVariant} tracks the LLM variants
available to the tool-calling orchestrator (singleton constraint: only one
active at a time).  A \textbf{ModelStats} entity stores per-model
evaluation snapshots, while a \textbf{DriftReport} entity records the
outcome of periodic distribution checks—creating a time-indexed
observability series for each model.

\paragraph{Dashboard cache.}

A \textbf{ModelResponse} entity caches the last prediction payload for
each combination of region, blood product, time horizon and alert level.
This avoids redundant inference calls and lets the dashboard render
instantly.

\paragraph{Data versioning.}

A lightweight \textbf{DataVersion} entity assigns a timestamp and a
content hash to each logical data source, giving downstream consumers a
cheap staleness check before re-fetching.

\paragraph{Blood-supply domain (file-based).}

Outside the relational store, three analytical entities model the supply
chain: \textbf{Donor} (demographics, eligibility, logistics),
\textbf{Hospital} (daily inventory flow, environmental context) and
\textbf{SupplyPoint} (stock, lead time, stockout indicators).  These live
as columnar Parquet files and are exposed through the feature store
(Section~\ref{sec:feature-store}).

\bigskip

Table~\ref{tab:mcd-relationships} summarises the inter-entity
relationships.

\begin{table}[htbp]
\centering
\caption{Entity relationships in the MCD.}
\label{tab:mcd-relationships}
\small
\begin{tabular}{@{}llcl@{}}
\toprule
\textbf{From} & \textbf{To} & \textbf{Card.} & \textbf{Design note} \\
\midrule
User            & Group / Permission   & $M{:}N$ & role-based access \\
Conversation    & User                 & $M{:}N$ & flexible thread membership \\
Message         & Conversation, User   & $N{:}1$ & immutable once created \\
AlertEvent      & AlertRule            & $N{:}1$ & nullable (manual alerts) \\
AlertEvent      & Notification         & $N{:}1$ & soft UUID link (decoupled) \\
AlertEventHistory & AlertEvent         & $N{:}1$ & append-only audit log \\
PredictionLog   & User, MLModelConfig  & $N{:}1$ & nullable (anonymous calls) \\
\bottomrule
\end{tabular}
\end{table}


% ─────────────────────────────────────────────────────────────────────────────
% 3.4.2  Modèle Logique (MLD)
% ─────────────────────────────────────────────────────────────────────────────

\subsection{Modèle Logique de Données (MLD)}
\label{sec:mld}

The MCD translates into two complementary storage tiers
(Figure~\ref{fig:mld-tiers}):

\begin{enumerate}
\item A \textbf{relational tier} (PostgreSQL in production, SQLite for
      development) managed by Django ORM, hosting all operational tables.
\item A \textbf{columnar tier} (Parquet files) holding the analytical
      datasets, accessed exclusively through the Feast feature store.
\end{enumerate}

\begin{figure}[htbp]
\centering
\includegraphics[width=\linewidth]{figures/mld_tiers_diagram.png}
\caption{Logical storage architecture.  Operational tables live in the
         relational database; analytical datasets live as Parquet files
         served through Feast.  The ML-backend microservice mirrors the
         model-registry tables so it can operate independently of
         Django.}
\label{fig:mld-tiers}
\end{figure}

This separation is deliberate.  Operational data (alerts, prediction logs,
model configs) requires transactional integrity, row-level updates and
foreign-key enforcement—properties that a relational database provides
naturally.  Analytical data (daily hospital inventories, donor snapshots,
stockout features) is append-heavy, read in bulk for training, and
benefits from columnar compression and predicate push-down—exactly what
Parquet delivers.

\paragraph{Relational tier — design choices.}

The backend is organised into seven Django applications, each owning a
cohesive set of tables.  Table~\ref{tab:mld-summary} lists every table
with its primary key strategy, the number of columns and the key design
decisions behind it.

\begin{table}[htbp]
\centering
\caption{Relational tables — summary.  PK strategy: \texttt{SERIAL}
         (auto-increment integer) or \texttt{UUID} (v4).  ``JSONB cols''
         counts columns stored as structured JSON.}
\label{tab:mld-summary}
\small
\begin{tabular}{@{}llccl@{}}
\toprule
\textbf{Table} & \textbf{App}
    & \textbf{PK} & \textbf{JSONB}
    & \textbf{Design rationale} \\
\midrule
\texttt{user}
    & auth & SERIAL & 0
    & email as USERNAME; custom manager \\
\texttt{conversation}
    & comms & SERIAL & 0
    & M2M join table for participants \\
\texttt{message}
    & comms & SERIAL & 0
    & FK to conversation + sender \\
\texttt{alert\_rule}
    & alerts & SERIAL & 2
    & conditions \& channels as JSONB \\
\texttt{alert\_event}
    & alerts & UUID & 1
    & UUID PK for cross-service refs \\
\texttt{alert\_event\_history}
    & alerts & SERIAL & 1
    & append-only; snapshot as JSONB \\
\texttt{notification}
    & notif & UUID & 1
    & UUID PK; decoupled from alerts \\
\texttt{model\_response}
    & dashboard & SERIAL & 1
    & composite index on scope dims \\
\texttt{ml\_model\_config}
    & ml & SERIAL & 4
    & features, info, examples, defaults \\
\texttt{prediction\_log}
    & ml & SERIAL & 2
    & input + output as JSONB \\
\texttt{orchestrator\_variant}
    & ml & SERIAL & 1
    & singleton constraint on selection \\
\texttt{model\_stats}
    & ml & SERIAL & 1
    & evaluation snapshot per model \\
\texttt{drift\_report}
    & ml & SERIAL & 1
    & per-feature detail as JSONB \\
\texttt{data\_version}
    & etl & SERIAL & 0
    & content hash for staleness check \\
\bottomrule
\end{tabular}
\end{table}

Several cross-cutting design decisions are worth noting:

\begin{itemize}
\item \textbf{UUID primary keys} are used for AlertEvent and Notification
      because these entities are referenced across service boundaries
      (Django $\leftrightarrow$ FastAPI); UUIDs avoid sequence collisions
      and make soft references safe.
\item \textbf{JSONB columns} are used wherever the schema is intentionally
      open-ended (alert conditions, feature metadata, prediction
      payloads).  This avoids an explosion of nullable columns while
      keeping the data queryable with JSON path operators.
\item \textbf{Composite indexes} are placed on the most common access
      patterns—e.g.\ \texttt{(severity, status, opened\_at)} on
      alert events, \texttt{(model\_id, checked\_at)} on drift reports,
      and \texttt{(region, product, period, alert\_level)} on dashboard
      responses.
\item \textbf{Soft deletion is avoided}; alert events transition through
      explicit status values, and prediction logs are append-only.  This
      simplifies queries and preserves a complete audit trail.
\end{itemize}

\paragraph{Registry mirror.}

The ML-backend microservice maintains a mirror of the model-registry
tables so it can serve predictions without depending on the Django
process.  When PostgreSQL is used, a database trigger pushes change
notifications to all connected workers in real time; when SQLite is used,
workers poll on a configurable interval.  This dual-write pattern provides
operational independence while keeping both services consistent.

\paragraph{Columnar tier.}

The analytical datasets are organised around three domain entities (Donor,
Hospital, SupplyPoint), each identified by a natural key plus an event
timestamp.  Data is stored in Parquet format, partitioned by entity type.
Table~\ref{tab:datasets-summary} lists the principal datasets.

\begin{table}[htbp]
\centering
\caption{Analytical datasets (Parquet).  Each row is a timestamped
         observation for one entity instance.}
\label{tab:datasets-summary}
\small
\begin{tabular}{@{}llrl@{}}
\toprule
\textbf{Dataset} & \textbf{Entity key}
    & \textbf{Cols} & \textbf{Prediction target(s)} \\
\midrule
Blood-bank daily ops
    & hospital + date & 23
    & stockout next day, heavy demand \\
Transfusion with features
    & hospital + date & 51
    & days until stockout, heavy demand \\
Donor registry
    & donor            & 27
    & donation within 6\,m, donation count \\
Donor logistics
    & donor            & 12
    & (joined to registry) \\
Donor snapshots
    & donor            & 9
    & donation within 30\,d \\
Hospital supply (ML-ready)
    & hospital + date  & 9
    & blood shortage (binary) \\
Stockout (ML-ready)
    & supply + date    & 9
    & days until stockout \\
\bottomrule
\end{tabular}
\end{table}


% ─────────────────────────────────────────────────────────────────────────────
% 3.4.3  Feature Store (offline / online)
% ─────────────────────────────────────────────────────────────────────────────

\subsection{Feature Store (offline / online)}
\label{sec:feature-store}

A persistent source of model-quality degradation in production is
\emph{training--serving skew}: features computed one way for training and a
slightly different way at inference time~\cite{sculley2015hidden}.  The
PIOS architecture eliminates this risk by introducing a \textbf{feature
store} (built on Feast) that acts as the single source of truth for every
feature consumed by any model.

\paragraph{Design goals.}

\begin{enumerate}
\item \textbf{Single definition.} Each feature is defined exactly once;
      both the offline (batch) and online (real-time) paths read from that
      definition.
\item \textbf{Point-in-time correctness.} Historical joins must respect
      causality—no future information may leak into a training row.
\item \textbf{Low-latency online serving.} Inference requests must
      retrieve the latest feature vector in milliseconds, not seconds.
\item \textbf{Extensibility.} Adding a new model that uses existing
      features should require zero feature re-engineering.
\end{enumerate}

\paragraph{Entity model.}

Three domain entities serve as join keys
(Table~\ref{tab:feast-entities}).  Every feature observation is anchored
to one entity instance and one event timestamp.

\begin{table}[htbp]
\centering
\caption{Feature-store entities.}
\label{tab:feast-entities}
\small
\begin{tabular}{@{}lll@{}}
\toprule
\textbf{Entity} & \textbf{Join key} & \textbf{Domain meaning} \\
\midrule
Donor      & \texttt{donor\_id}    & Individual blood donor \\
Hospital   & \texttt{hospital\_id} & Blood-bank / hospital site \\
Supply     & \texttt{supply\_id}   & Inventory supply point \\
\bottomrule
\end{tabular}
\end{table}

\paragraph{Feature views and services.}

For each entity one or more \emph{feature views} expose a typed schema
over a Parquet source file.  Every view carries a time-to-live (TTL) that
controls how long a feature value remains valid for online serving.
Views sharing the same entity are bundled into a \emph{feature service}—a
logical group that training and inference code can reference by name.
Three services are defined:

\begin{itemize}
\item \textbf{Donor service} — demographics, geolocation, donation
      recency, frequency and cluster assignment;
\item \textbf{Hospital service} — environmental conditions (temperature,
      rain), demand drivers (surgeries, trauma), current inventory;
\item \textbf{Supply service} — stock level, daily usage, lead time,
      restock recency, upcoming surgical demand.
\end{itemize}

This grouping means a new model—say, a donor-eligibility classifier—can
reference the Donor service and immediately receive all donor features
without any schema duplication.

\paragraph{Offline path.}

During training a script provides an entity dataframe with join keys and
label timestamps.  The feature store performs a \emph{point-in-time join}:
for each row it retrieves the most recent feature snapshot whose timestamp
is \emph{strictly before} the label's timestamp.  This guarantees that the
model is trained only on information that would have been available at
decision time.

\paragraph{Online path.}

For real-time inference the feature store maintains a materialised online
store.  When a prediction request arrives, the serving layer retrieves the
latest feature vector for the given entity key with sub-millisecond
latency.  Because the feature definitions are shared with the offline
path, training--serving skew is eliminated by construction.

\paragraph{Materialisation strategy.}

Features are pushed from the offline source to the online store via
incremental materialisation: only rows newer than the previous
high-water-mark timestamp are transferred.  This makes the refresh cost
proportional to the volume of new data rather than the total dataset size,
and it is triggered automatically before every training run so the online
store is never stale when inference begins.


% ─────────────────────────────────────────────────────────────────────────────
% 3.4.4  Gestion des données temporelles
% ─────────────────────────────────────────────────────────────────────────────

\subsection{Gestion des données temporelles}
\label{sec:temporal-management}

Blood-supply data is intrinsically temporal: donations occur at specific
instants, inventory levels change hour by hour, and the relevance of a
prediction decays rapidly.  The PIOS data architecture addresses this
through four complementary design decisions.

\paragraph{1.\ Event-time anchoring.}

Every feature record carries an event timestamp that anchors it to a
specific moment in the real world.  The feature store uses this column for
point-in-time joins (Section~\ref{sec:feature-store}), and the TTL
mechanism uses it to expire stale values from the online store.  Together,
these two mechanisms prevent both future leakage (at training time) and
stale-data pollution (at inference time).

\paragraph{2.\ Causal feature engineering.}

The raw time-series datasets (daily hospital operations, donor snapshots)
are enriched with a family of temporal features designed to capture
short-term dynamics and longer-term trends while respecting causality
(Table~\ref{tab:temporal-features}).  All lag and rolling computations are
performed per entity and sorted by date, so no future information leaks
into past rows.

\begin{table}[htbp]
\centering
\caption{Families of temporal features.  All windows are strictly causal.}
\label{tab:temporal-features}
\small
\begin{tabular}{@{}lp{7.2cm}@{}}
\toprule
\textbf{Family} & \textbf{Purpose \& examples} \\
\midrule
Calendar
    & Encode seasonality: day-of-week, weekend, month, holiday \\
Lag (1 / 7 / 30\,d)
    & Capture recent state: yesterday's usage, last week's stock \\
Rolling mean \& std
    & Smooth noise and detect volatility over 7- and 30-day windows \\
Trend (OLS slope)
    & Distinguish rising from falling demand (7-day, 30-day) \\
Coverage ratio
    & Days of supply remaining at current consumption rate \\
Demand pressure
    & Surgeries and trauma relative to available stock \\
Stockout memory
    & Recent stockout flags and 7-day cumulative count \\
Net flow
    & Units collected minus units used (single-step balance) \\
\bottomrule
\end{tabular}
\end{table}

\paragraph{3.\ Scheduled retraining and drift monitoring.}

As the real-world data distribution evolves (\emph{concept drift}),
models must be retrained periodically.  The platform includes a background
scheduler that supports daily, weekly and monthly retraining cadences.
After each training run a four-step pipeline executes:

\begin{enumerate}
\item \textbf{Train} — produce a new model registered as
      \emph{challenger};
\item \textbf{Promote} — compare the challenger against the current
      \emph{champion} on a primary metric; promote only if the
      challenger is at least as good;
\item \textbf{Sync} — propagate the new champion to the inference
      registry so production traffic uses it immediately;
\item \textbf{Drift check} — compare feature distributions between
      the reference (training-time) and current (production) data
      using four statistical tests
      (Table~\ref{tab:drift-tests}).
\end{enumerate}

\begin{table}[htbp]
\centering
\caption{Statistical tests for drift detection.}
\label{tab:drift-tests}
\small
\begin{tabular}{@{}lll@{}}
\toprule
\textbf{Test} & \textbf{Feature type} & \textbf{Alert threshold} \\
\midrule
Kolmogorov--Smirnov    & numerical    & $p < 0.05$ \\
Population Stability Index & numerical & PSI $> 0.20$ \\
Chi-squared            & categorical  & $p < 0.05$ \\
Jensen--Shannon div.   & categorical  & JS $> 0.10$ \\
\bottomrule
\end{tabular}
\end{table}

Each feature is assigned a severity (\emph{none}, \emph{warning},
\emph{critical}).  When critical drift is detected the alert subsystem
(Section~\ref{sec:mcd}) can trigger immediate retraining outside the
regular schedule.

\paragraph{4.\ Simulator temporal validation.}

The discrete-event simulator produces 6-hourly inventory snapshots.
Before trusting any simulation-based evaluation, these series are
validated for temporal plausibility using (a)~an Augmented Dickey--Fuller
test for stationarity—confirming that stock levels do not drift
unboundedly—and (b)~Ljung--Box / ACF analysis to verify realistic
short-range autocorrelation at 6-hour multiples.  A 72-hour warm-up is
discarded to remove the initial transient.

\bigskip

Together, event-time anchoring, causal feature engineering, scheduled
drift monitoring and simulator temporal validation form a coherent design
that ensures every component of the platform—from training through
inference to simulation—handles time-dependent data correctly.
